\subsection{Event-driven Selective Semantic Updates}
\label{sec:semantic_physical_coordination}

Task-relevant events enter the shared semantic state, whereas continuous
physical evolution remains within manipulator-local execution. Event handling
identifies the invalid scope and invokes adaptation, as shown in
Fig.~\ref{fig:event_graph_update}.

\begin{figure}[t!]
  \centering
  \definecolor{figink}{HTML}{20252A}
\definecolor{figmuted}{HTML}{6B737A}
\definecolor{figpurple}{HTML}{8E5BBE}
\definecolor{figgreen}{HTML}{4D8B55}
\definecolor{figblue}{HTML}{2F6DA6}
\definecolor{figline}{HTML}{D9DDE3}

% Fill one IEEE/ICRA column exactly.  The explicit reset of the TikZ
% bounding box removes the residual left margin introduced by node extents.
\resizebox{\columnwidth}{!}{%
\begin{tikzpicture}[x=1cm,y=1cm,>=Latex]
  \tikzset{
    head/.style={font=\sffamily\bfseries\fontsize{8}{8.8}\selectfont,text=figink},
    subtitle/.style={font=\sffamily\fontsize{8}{8.8}\selectfont,text=figink,align=center},
    eventlabel/.style={font=\sffamily\bfseries\fontsize{8}{8.8}\selectfont,
                       text=figink},
    flow/.style={-Latex,draw=figink!70,line width=.42pt},
    graphpdf/.style={inner sep=0pt,outer sep=0pt}
  }

  % Six graph PDFs are stored under figures/figure3/.
  \def\figpath{figures/figure3}

  % ------------------------------------------------------------
  % Column headers
  % ------------------------------------------------------------
  \node[head] at (1.85,6.82) {$G$};
  \node[head] at (6.94,6.82) {$G'$};

  % ------------------------------------------------------------
  % E1: Task update
  % ------------------------------------------------------------
  \node[graphpdf] at (1.85,5.58)
    {\includegraphics[width=3.70cm]{\figpath/figure3_E1_Gassign_before.pdf}};
  \draw[flow] (3.73,5.58) -- node[eventlabel,above=2pt] {$e^{\texttt{t}}$} (5.00,5.58);
  \node[graphpdf] at (6.94,5.58)
    {\includegraphics[width=3.84cm]{\figpath/figure3_E1_Gassign_after.pdf}};
  \node[subtitle,text width=8.56cm] at (4.43,4.50)
    {\textbf{Task update:} the user wants crisps instead of biscuit after $n_4$ is executed.};

  \draw[figline,line width=.35pt] (0.00,4.15) -- (8.86,4.15);

  % ------------------------------------------------------------
  % E2: External disturbance
  % ------------------------------------------------------------
  \node[graphpdf] at (1.85,3.19)
    {\includegraphics[width=3.70cm]{\figpath/figure3_E2_Gassign_before.pdf}};
  \draw[flow] (3.73,3.19) -- node[eventlabel,above=2pt] {$e^{\texttt{x}}$} (5.00,3.19);
  \node[graphpdf] at (6.94,3.19)
    {\includegraphics[width=3.84cm]{\figpath/figure3_E2_Gassign_after.pdf}};
  \node[subtitle,text width=8.56cm] at (4.43,2.06)
    {\textbf{External disturbance:} the biscuit is manually taken and placed downstream on the conveyor while $n_5$ is being executed.};

  \draw[figline,line width=.35pt] (0.00,1.72) -- (8.86,1.72);

  % ------------------------------------------------------------
  % E3: Execution failure
  % ------------------------------------------------------------
  \node[graphpdf] at (1.86,0.67)
    {\includegraphics[width=3.72cm]{\figpath/figure3_E3_Gassign_before.pdf}};
  \draw[flow] (3.75,0.67) -- node[eventlabel,above=2pt] {$e^{\texttt{f}}$} (4.98,0.67);
  \node[graphpdf] at (6.93,0.67)
    {\includegraphics[width=3.86cm]{\figpath/figure3_E3_Gassign_after.pdf}};
  \node[subtitle,text width=8.60cm] at (4.43,-0.63)
    {\textbf{Execution failure:} $n_1$ fails because the conveyor moves too fast and the Franka misses the grasp; the downstream suction arm prepares in advance for the pick.};

  % Reset the natural node-based bounding box, then define an exact full-column box.
  % This removes the visible blank strip on the left when the input is centered.
  \pgfresetboundingbox
  \path[use as bounding box] (0.00,-0.90) rectangle (8.86,7.02);
\end{tikzpicture}%
}

  \caption{Selective graph transitions
  $G\rightarrow e_k\rightarrow G'$ for task updates, external
  disturbances, and execution failures. Purple contours mark
  $\mathcal{N}_k^{\texttt{aff}}$, unfinished or dependency-invalidated nodes
  requiring reconsideration and reconstruction; failed attempts recorded in
  $z$ need not be reconstructed. Green nodes are retained and blue nodes are new.}
  \label{fig:event_graph_update}
\end{figure}

\subsubsection{Event-affected Scope Identification}
\label{sec:affected_subtasks}

The event types are defined in~\eqref{eq:semantic_events}. A task update
changes the current instruction, an external disturbance is interpreted from
fresh visual evidence, and an execution failure is established by verification.
The event then enters the verification role, which identifies the affected
assignment-aware primitive nodes:
\begin{equation}
 \mathcal{N}_k^{\texttt{aff}}\triangleq
 \rho_V^{\texttt{aff}}\bigl(
 e_k,\gamma,G,\mathcal{F},z\bigr),
 \label{eq:affected_subtasks}
\end{equation}
where $\mathcal{N}_k^{\texttt{aff}}\subseteq\mathcal{N}$
contains the nodes whose task meaning, feasibility, execution status, or
dependencies require reconsideration. The operator tests $e_k$ against the
current $\gamma$, $G$, refreshed $\mathcal{F}$, and verified $z$. Its
event-specific inconsistency checks select requirements invalidated by a task
update, scene-dependent assumptions invalidated by a disturbance, or nodes
marked $\texttt{FAILED}$ after execution. Verification then traces outgoing
dependencies in $G$ to include unfinished descendants whose validity or
executability depends on those nodes, while independent valid nodes remain
outside $\mathcal{N}_k^{\texttt{aff}}$ and are preserved.

\subsubsection{Selective Semantic-state Reconstruction}
\label{sec:update_affected_state}

The selective update produces the next committed semantic state:
\begin{equation}
 S_{\nu+1}^{\texttt{sem}}\triangleq
 \mathcal{U}\!\left(
 S_\nu^{\texttt{sem}},e_k,\mathcal{N}_k^{\texttt{aff}}\right),
 \label{eq:semantic_state_update}
\end{equation}
where $\mathcal{U}$ denotes the final result of the coordinated role updates
rather than a separate replanning step. Before this result is committed, the
affordance role refreshes $\mathcal{F}$, and the verification role updates $z$
and produces $\mathcal{N}_k^{\texttt{aff}}$ by~\eqref{eq:affected_subtasks}.
Using this affected set and the latest $\mathcal{F}$, the planning role retains
exactly the nodes of $G$ outside $\mathcal{N}_k^{\texttt{aff}}$, together with
the relations whose endpoints are both retained. It replans the nodes in
$\mathcal{N}_k^{\texttt{aff}}$, grounds and capability-assigns their
replacements from the updated state, and reconnects the revised subgraph to the
retained graph at the dependency boundary, thereby completing a partial
reconstruction of $G$. The status of each retained node is preserved, whereas
each new or revised node starts as $\texttt{PENDING}$. The reconstructed graph
is checked for capability-valid assignments and predecessor consistency. The
affordance, planning, and verification roles have then respectively updated
$\mathcal{F}$, $G$, and $z$; together with the event-consistent $\gamma$, these
components form $S_{\nu+1}^{\texttt{sem}}$ and advance the semantic version from
$\nu$ to $\nu+1$.

\subsection{Primitive-command Dispatch and Execution}
\label{sec:dispatch_execute}

\subsubsection{Capability-valid Command Construction}
\label{sec:construct_commands}

The command notion  is
instantiated after planning has produced assignment-aware primitive nodes. At
each execution step, $\mathcal{N}_{\texttt{exec}}$ is the temporary set of
executable primitive subtasks in $G$. A node $n$
belongs to $\mathcal{N}_{\texttt{exec}}$ when
$z(n)\ne\texttt{COMPLETE}$ and $z(n')=\texttt{COMPLETE}$ for every
$n'\in\operatorname{pred}(n)$. For any
$n\in\mathcal{N}_{\texttt{exec}}$, command construction is given by:
\begin{equation}
 c_n=\phi_C(n),
 \qquad
 \phi_C:\mathcal{N}_{\texttt{exec}}\rightarrow\mathcal{C},
 \label{eq:execution_command}
\end{equation}
where $\phi_C$ is the command construction function and $\mathcal{C}$ denotes
the executable command space, distinct from the per-manipulator capability
sets $\{\mathcal{C}_i\}_{i=1}^{N}$. Although $n$ contains semantic execution
information, $c_n$ is the executable command sent to the robot interface. At
each execution step, the commands generated from
$\mathcal{N}_{\texttt{exec}}$ form a concurrent command set
$\mathbf{c}_m=\{c_{n_1},c_{n_2},\cdots\}$.

\subsubsection{Local Execution and Outcome Verification}
\label{sec:manipulator_local_execution}

The assigned local skill executes each command using the manipulator's current
perception, tracking, and control stack. Its target reference allows continuous
motion to remain local. After execution, the latest observation refreshes
$\mathcal{F}$ through $\mathcal{F}\gets\rho_A(o_t')$ without semantic
replanning. For each executed command, verification extracts $\eta_n$ via
\eqref{eq:verification_evidence} and uses $(\eta_n,\gamma)$ to update $z(n)$ and
detect any $e_k$ via~\eqref{eq:verification_status}. A detected event triggers
identification and selective state revision.

\subsection{Closed-Loop Coordination}
\label{sec:closed_loop_coordination}

\subsubsection{Online Coordination Algorithm}

Algorithm~\ref{alg:coordination} summarizes the online procedure. At task start,
the affordance role observes the scene and extracts $\mathcal{F}$. The planning
role then derives the task dependencies, constructs the assignment-aware $G$,
and initializes $S_\nu^{\texttt{sem}}=(\gamma,G,\mathcal{F},z)$. During each
loop, the verified frontier selects independent executable nodes, maps them to
capability-valid commands, and dispatches the concurrent command set to the
assigned manipulators. After local execution, affordance refreshes
$\mathcal{F}$; for every executed command, verification extracts $\eta_n$ and
uses $(\eta_n,\gamma)$ to update $z(n)$ and detect any $e_k$. When an event
occurs, verification reads the refreshed $\mathcal{F}$ together with
$e_k$, $\gamma$, $G$, and $z$, applies the event-specific criterion and traces
unfinished dependent nodes, and outputs $\mathcal{N}_k^{\texttt{aff}}$. The
planning role then retains the nodes outside this set, replans only the affected
nodes using the latest $\mathcal{F}$, and reconnects the revised subgraph. The
three role updates are committed as $S_{\nu+1}^{\texttt{sem}}$, after which
dispatch resumes from the new consistent frontier. Restricting recomputation to this
dependency-closed scope is intended to shorten semantic-event response time and
reduce the chance that ongoing physical evolution removes an available
manipulation opportunity.

\begin{algorithm}[t!]
\caption{Online role-shared semantic coord.}
\label{alg:coordination}
\KwIn{$\gamma$, $o(t)$, $\{\mathcal{C}_i\}_{i=1}^{N}$, $\mathcal{R}$}
\KwOut{$\mathbf{C}$}

Extract $\mathcal{F}$ from $o(0)$\;
Construct $G$ by~\eqref{eq:planning_pipeline}\;
Initialize $S_\nu^{\texttt{sem}}$ by~\eqref{eq:semantic_state};
$\mathbf{C}\gets()$\;

\While{task incomplete}{
    $\mathcal{N}_{\texttt{exec}}\gets
    \{n\in\mathcal{N}\mid z(n)\ne\texttt{COMPLETE},
    \ z(n')=\texttt{COMPLETE},\
    \forall n'\in\operatorname{pred}(n)\}$\;

    $\mathbf{c}_m\gets
    \{c_n\mid n\in\mathcal{N}_{\texttt{exec}}\}$
    using~\eqref{eq:execution_command}\;
    Append $\mathbf{c}_m$ to $\mathbf{C}$ and execute it\;
    Refresh $\mathcal{F}$ from $o_t'$\;

    \ForEach{$c_n\in\mathbf{c}_m$}{
        Extract $\eta_n$ by~\eqref{eq:verification_evidence}\;
        Update $(e_k,z(n))$ by~\eqref{eq:verification_status}\;
    }

    \If{$e_k$ is detected}{
        Identify $\mathcal{N}_k^{\texttt{aff}}$
        by~\eqref{eq:affected_subtasks}\;
        Retain $G\setminus\mathcal{N}_k^{\texttt{aff}}$ and reconstruct
        the affected subgraph using $\mathcal{F}$\;
        Commit $S_{\nu+1}^{\texttt{sem}}$
        by~\eqref{eq:semantic_state_update};
        $\nu\gets\nu+1$\;
    }
}
\end{algorithm}

\subsubsection{Event-specific Coordination}
\label{sec:event_specific_algorithm}

The same protocol handles all events. For $\texttt{E1}$, verification marks the
requirement-changing branch and planning regenerates only that branch. For
$\texttt{E2}$, refreshed $\mathcal{F}$ invalidates nodes whose target,
feasibility, or dependency assumptions changed. For $\texttt{E3}$, the failed
node and unfinished downstream descendants receive a  repair
branch. Valid completed nodes, independent pending work, and local tracking are
retained. Thus, affordance supplies current facts, verification selects scope,
and planning revises only the invalidated semantic subgraph.

\subsubsection{Computational Complexity}

The explicit graph operations in Algorithm~\ref{alg:coordination} are linear
in the maintained graph when predecessor lists are stored with the graph.
Checking executable nodes requires scanning $\mathcal{N}$ and the relevant
incoming dependencies, with worst-case cost $O(|\mathcal{N}|+|\mathcal{L}|)$ per
dispatch step. After an event, selective revision is bounded by the affected
subgraph identified by $\mathcal{N}_k^{\texttt{aff}}$ and its incident
dependencies. The complexity of the semantic operators
$\rho_P$, $\rho_V^{\texttt{e}}$,
$\rho_V^{\texttt{z}}$, and $\rho_V^{\texttt{aff}}$ also includes the call time
of the concrete LLM/VLM models used by these operators. Because these call
times depend on the selected models and serving conditions, no closed-form
bound is assigned here.
